\documentclass[11pt]{article}

\usepackage[margin=1in]{geometry}
\usepackage{amsmath, amssymb, amsthm}
\usepackage{mathtools}
\usepackage{bm}
\usepackage{booktabs}
\usepackage{array}
\usepackage{enumitem}

\DeclareMathOperator*{\clip}{clip}
\DeclareMathOperator{\sign}{sgn}
\newcommand{\Ehat}{\widehat{\mathbb{E}}}
\newcommand{\dbudget}{\hat d_{\text{budget}}}

\title{BSPO --- Iter2 Formal Algorithm Specification\\Single-Domain}
\author{}
\date{}

\begin{document}
\maketitle

% =============================================================================
\section{Iter2 changes from iter1}

Iter2 is a refinement of the single-domain BSPO spec based on iter1
empirical observations on the B300 dev cluster (see
\texttt{bisimpo/sd\_ablation\_report.tex} for iter1 data). Changes:

\begin{enumerate}[label=(\alph*)]
\item \textbf{Drop $\Delta_\tau^{(\mathrm W)}$.} Iter1 exposed three
  variants (simplest, w\_penalty using $\mathrm{sgn}(s_\tau)\cdot
  \clip(\min(s_\tau^+,s_\tau^-),\delta)$, and w\_penalty\_only using
  $s_\tau\cdot\hat A_\tau$). Iter2 keeps only the regular clip
  $\Delta_\tau^{(\mathrm{reg})}$. The ``w\_penalty'' iter2 variant now
  means $\Delta_\tau^{(\mathrm{reg})}\hat A_\tau - \lambda\cdot\mathrm{pen}_\tau$
  (not the iter1 sign-min form).
\item \textbf{Per-group and per-trajectory clip thresholds.} Iter1 used
  a single flat $\delta$ for all levels. Iter2 introduces
  $\epsilon_\tau$ and $\epsilon_g$ derived from hyperparameters
  $\delta_{\rm Tj}, \delta_{\rm Grp}$ and a budget allocation
  $\dbudget$ over groups.
\item \textbf{$\mathrm{pen}_\tau$ uses $\epsilon_{g(\tau)}$}, not
  $\epsilon_\tau$. Intentional: the penalty activates when a
  trajectory's deviation exceeds its \emph{group} budget, not its own
  trajectory budget.
\item \textbf{Four advantage variants.} Iter1 used GRPO standard only
  (Adv~I below). Iter2 adds three alternatives (Adv~II/III/IV) that
  replace the standard-deviation denominator with a reward-spread
  measure $R_{\max\text{res}}(dom)$ and/or the group dispersion
  $\hat d(g)$. The set of legal advantages depends on the chosen
  budget allocation (pro/const/ipro).
\item \textbf{Sync / async clip} toggle. Sync (default): a single
  $\epsilon_\tau$ applied symmetrically on both sides of the clip.
  Async: $\epsilon_\tau$ is chosen by the sign of $\hat A_\tau$, to
  match the asymmetric GSPO reference threshold
  ($3\!\times\!10^{-4}$ for positive advantage,
  $4\!\times\!10^{-4}$ for negative).
\item \textbf{Two labeled variants} for iter2: $V_{1a}$ (simplest) and
  $V_{1b}$ (w\_penalty). The $V_{1a/b}$ labels are within iter2's
  objective space. The $V_1\ldots V_5$ labels from iter1 retain
  their iter1 meaning and are not reused.
\end{enumerate}

% =============================================================================
\section{Single-domain BSPO (iter2)}

% -----------------------------------------------------------------------------
\subsection{Inputs and notation}

\begin{itemize}[noitemsep]
\item Batch of $B$ trajectories. Each trajectory $\tau$ has token count
  $|\tau|$, token-level policy ratio
  $r_t = \pi_\theta(a_t\mid s_t)/\pi_{\theta_{\text{old}}}(a_t\mid s_t)$,
  and scalar reward $R_\tau$.
\item Groups: trajectories are partitioned by prompt id; $g(\tau)$
  denotes the group of $\tau$. In single-domain, the sole domain is
  denoted $dom$; $dom(\tau)=dom(g)=dom$ for every $\tau,g$.
\item Token counts: $T_g = \sum_{\tau\in g}|\tau|$,
  $T_{dom}=\sum_{g\in dom}T_g$.
\end{itemize}

% -----------------------------------------------------------------------------
\subsection{Ratio statistics}

\begin{align*}
s_\tau^{+} &= \frac{1}{|\tau|}\sum_{t=1}^{|\tau|}\max(r_t-1,\,0),
& s_\tau^{-} &= \frac{1}{|\tau|}\sum_{t=1}^{|\tau|}\max(1-r_t,\,0),
& s_\tau &= s_\tau^{+} - s_\tau^{-}.
\end{align*}

% -----------------------------------------------------------------------------
\subsection{Budget allocation}

\paragraph{Per-level dispersion.}
\[
\hat d(z) \;=\; \frac{1}{|z|}\sum_{\tau\in z}\bigl|R_\tau - \bar R_z\bigr|,
\qquad z\in\{g, dom\}.
\]

\paragraph{Per-group token weight within a domain.}
\[
\hat w(g/dom) \;=\; \frac{T_g^{2}}{T_{dom}^{2}}\,\hat d(g).
\]

\paragraph{Per-domain budget.}
\[
\dbudget(dom) \;=\; \frac{\hat d(dom)}{R_{\max\text{res}}(dom)},
\qquad
R_{\max\text{res}}(dom) \;=\; \max_{\tau\in dom}\bigl|R_\tau - \bar R_{dom}\bigr|.
\]

\paragraph{Per-group budget --- three allocation variants.}
The allocation redistributes the domain budget across groups in one of
three shapes.

\begin{equation}
\dbudget(g) \;=\;
\begin{dcases}
\dfrac{\hat w(g/dom)}{\sum_{g'\in dom}\hat w(g'/dom)^{2}}
  & \text{\textbf{pro} --- proportional to $\hat w(g/dom)$} \\[6pt]
\dfrac{1}{\sum_{g'\in dom}\hat w(g'/dom)}
  & \text{\textbf{const} --- equal across groups} \\[6pt]
\dfrac{1}{T_g\,\hat w(g/dom)}
  & \text{\textbf{ipro} --- inverse proportional to $T_g\,\hat w(g/dom)$}
\end{dcases}
\label{eq:dbudget-g}
\end{equation}

% -----------------------------------------------------------------------------
\subsection{Clip thresholds}

\paragraph{Trajectory and group thresholds.}
\begin{align}
\epsilon_\tau &= \dbudget(dom(\tau))\;\dbudget(g(\tau))\;\delta_{\rm Tj}
                    \;\sim\; \epsilon_{\rm GSPO}
                    \label{eq:eps-tau} \\
\epsilon_g    &= \dbudget(dom(g))\;\dbudget(g)\;\delta_{\rm Grp}
                    \;\sim\; 3\!\times\!10^{-3}.
                    \label{eq:eps-g}
\end{align}

The $\sim$ is a design guidance: the hyperparameters $\delta_{\rm Tj}$
and $\delta_{\rm Grp}$ are \emph{not} fixed constants; they are tuned
so that the resulting $\epsilon_\tau$ and $\epsilon_g$ sit at the
target scales. The target scales are:
\begin{itemize}[noitemsep]
\item $\epsilon_\tau$ near the GSPO reference clip $\epsilon_{\rm GSPO}$
  ($3\!\times\!10^{-4}$ for positive advantage,
  $4\!\times\!10^{-4}$ for negative, see \S\ref{sec:async-clip}).
\item $\epsilon_g$ roughly an order of magnitude larger, near
  $3\!\times\!10^{-3}$.
\end{itemize}

\paragraph{Practical tuning notes.}
\begin{enumerate}[label=(\roman*)]
\item Given $\epsilon_{\rm GSPO} = 3\!\times\!10^{-4}$, once
  $\delta_{\rm Tj}$ is set so that $\epsilon_\tau$ lands at this value,
  $\delta_{\rm Grp}$ can be taken to be
  $10\!\times\!\delta_{\rm Tj}$ (and $\epsilon_g = 10\epsilon_\tau$
  follows because the $\dbudget$ factors cancel between \eqref{eq:eps-tau}
  and \eqref{eq:eps-g}).
\item The three $\dbudget(g)$ allocation cases
  (\textbf{pro}, \textbf{const}, \textbf{ipro}) produce different
  absolute magnitudes, so $\delta_{\rm Tj}$ and $\delta_{\rm Grp}$
  must be recomputed per allocation to keep $\epsilon_\tau,\epsilon_g$
  at their target scales.
\item The $\sim$ guidance in \eqref{eq:eps-tau}--\eqref{eq:eps-g}
  means: tweak $\delta_{\rm Tj},\delta_{\rm Grp}$ empirically so
  $\epsilon_\tau$ is on the order of $\epsilon_{\rm GSPO}$ and
  $\epsilon_g$ is on the order of $3\!\times\!10^{-3}$ --- an exact
  match is neither required nor guaranteed under all
  $(\hat d,\,\hat w)$ profiles.
\end{enumerate}

\paragraph{Sync vs async clip.}\label{sec:async-clip}
Iter2 supports two clipping modes. The asymmetric
$\epsilon_{\rm GSPO}$ ($3e\text{-}4$ pos, $4e\text{-}4$ neg) is
realised \emph{only} in async mode.

\begin{description}
\item[Sync (default):] a single $\delta_{\rm Tj}$ (and
  $\delta_{\rm Grp}$) produces one scalar $\epsilon_\tau$ (and
  $\epsilon_g$) per trajectory. Used for every trial unless
  \textsf{async\_clip} is toggled.
\item[Async:] two hyperparameter values per level:
  $\delta_{\rm Tj}^{(+)},\delta_{\rm Tj}^{(-)}$ and
  $\delta_{\rm Grp}^{(+)},\delta_{\rm Grp}^{(-)}$.
  For a trajectory $\tau$ with advantage $\hat A_\tau$:
  \begin{align*}
  \epsilon_\tau &= \dbudget(dom)\,\dbudget(g(\tau))\,
                   \delta_{\rm Tj}^{(\sign(\hat A_\tau))},\\
  \epsilon_{g(\tau)} &= \dbudget(dom)\,\dbudget(g(\tau))\,
                   \delta_{\rm Grp}^{(\sign(\hat A_\tau))}.
  \end{align*}
  Targets: $\epsilon_\tau \sim 3\!\times\!10^{-4}$ if
  $\hat A_\tau>0$, else $\sim 4\!\times\!10^{-4}$.
  $\epsilon_g$ 10$\times$ larger on both sides.
  Both sides of each single-trajectory clip (i.e.\ $s_\tau^+$ and
  $s_\tau^-$) share the same $\epsilon_\tau$; the asymmetry is
  \emph{between trajectories}, driven by the advantage sign.
\end{description}

% -----------------------------------------------------------------------------
\subsection{Advantage variants}

GRPO-style per-trajectory advantages. Four cases are supported; legal
pairings with the $\dbudget(g)$ allocation are listed in
Table~\ref{tab:adv-alloc}.

\begin{align}
\textbf{Adv~I (default)}\;:\;
\hat A_\tau &= \dfrac{R_\tau - \bar R_{g(\tau)}}{\sigma_{g(\tau)} + \varepsilon}
  \label{eq:adv-I} \\[2pt]
\textbf{Adv~II}\;:\;
\hat A_\tau &= \dfrac{R_\tau - \bar R_{g(\tau)}}
                    {\hat d(g(\tau))\,R_{\max\text{res}}(dom)}
  \label{eq:adv-II} \\[2pt]
\textbf{Adv~III}\;:\;
\hat A_\tau &= \dfrac{R_\tau - \bar R_{g(\tau)}}{R_{\max\text{res}}(dom)}
  \label{eq:adv-III} \\[2pt]
\textbf{Adv~IV}\;:\;
\hat A_\tau &= \hat d(g(\tau))\,\dfrac{R_\tau - \bar R_{g(\tau)}}{R_{\max\text{res}}(dom)}
  \label{eq:adv-IV}
\end{align}
\[
\text{where}\quad
\bar R_g = \tfrac{1}{|g|}\sum_{\tau\in g}R_\tau,\quad
\sigma_g^{2} = \tfrac{1}{|g|}\sum_{\tau\in g}(R_\tau - \bar R_g)^{2}.
\]

\begin{table}[h]
\centering
\caption{Legal advantage $\times$ allocation pairings.}
\label{tab:adv-alloc}
\begin{tabular}{lcccc}
\toprule
                 & Adv I & Adv II & Adv III & Adv IV \\
\midrule
\textbf{pro}    & \checkmark & \checkmark & \checkmark & \checkmark \\
\textbf{const}  & \checkmark & ---        & ---        & ---        \\
\textbf{ipro}   & \checkmark & \checkmark & ---        & ---        \\
\bottomrule
\end{tabular}
\end{table}

% -----------------------------------------------------------------------------
\subsection{Clipping operators}

\paragraph{Regular clip $\Delta_\tau^{(\mathrm{reg})}$.}
\begin{equation}
\Delta_\tau^{(\mathrm{reg})} \;=\; \clip(s_\tau^{+},\,\epsilon_\tau)
                               \;-\; \clip(s_\tau^{-},\,\epsilon_\tau).
\label{eq:delta-reg}
\end{equation}
In sync mode $\epsilon_\tau$ is one scalar per trajectory; in async
mode $\epsilon_\tau$ is drawn from $\{\epsilon_\tau^{(+)},
\epsilon_\tau^{(-)}\}$ by $\sign(\hat A_\tau)$. Both sides of the clip
in \eqref{eq:delta-reg} always use the same chosen $\epsilon_\tau$.

\paragraph{Penalty $\mathrm{pen}_\tau$.}
\begin{equation}
\mathrm{pen}_\tau \;=\; \max\!\left(0,\;
  \frac{\min(s_\tau^{+},s_\tau^{-})}{\epsilon_{g(\tau)}} - 1\right).
\label{eq:pen-tau}
\end{equation}

Note the denominator: $\epsilon_{g(\tau)}$, not $\epsilon_\tau$. This
is the deliberate iter2 change: the penalty fires when the bi-lateral
deviation $\min(s_\tau^+, s_\tau^-)$ exceeds the
\emph{group}-level budget. Because $\epsilon_g \approx 10\,\epsilon_\tau$,
most trajectories with normal fluctuations around $\epsilon_\tau$ will
\emph{not} trigger the penalty; only trajectories whose bi-lateral
deviation approaches the group budget do.

% -----------------------------------------------------------------------------
\subsection{Per-trajectory objective}

\begin{align}
V_{1a}\;\text{(simplest):}\quad
\ell_\tau &= \Delta_\tau^{(\mathrm{reg})}\,\hat A_\tau
  \label{eq:v1a} \\[2pt]
V_{1b}\;\text{(w\_penalty):}\quad
\ell_\tau &= \Delta_\tau^{(\mathrm{reg})}\,\hat A_\tau
            \;-\;\lambda\,\mathrm{pen}_\tau
  \label{eq:v1b}
\end{align}

$V_{1a}$ and $V_{1b}$ differ \emph{only} by the penalty term; both use
the same $\Delta^{(\mathrm{reg})}$ and the same chosen $\hat A_\tau$.
$V_{1b}$ reduces to $V_{1a}$ when $\lambda=0$.

% -----------------------------------------------------------------------------
\subsection{Aggregated loss}

\[
\hat J \;=\; \frac{1}{B}\sum_{i=1}^{B}\ell_{\tau_i},
\qquad
\mathcal L \;=\; -\hat J.
\]

% =============================================================================
\section{Hyperparameters and known constants}

\paragraph{Tunable.}
\[
\delta_{\rm Tj} > 0,\quad \delta_{\rm Grp} > 0,\quad \lambda \ge 0.
\]
In async mode these split into signed pairs:
$\delta_{\rm Tj}^{(+)}, \delta_{\rm Tj}^{(-)},
\delta_{\rm Grp}^{(+)}, \delta_{\rm Grp}^{(-)}$.

\paragraph{Known constants (single-domain math, iter2 trial configuration).}
\[
T_g = 16,\qquad
T_{dom} = 128\times 16 = 2048,\qquad
R_{\max\text{res}}(dom) = 1.
\]

$R_{\max\text{res}}(dom)$ should in general be computed during a
profiling pass; for the single-domain math setup it is given a priori
($=1$). The formula is retained for reference:
\[
R_{\max\text{res}}(dom) \;=\; \max_{\tau\in dom}\bigl|R_\tau - \bar R_{dom}\bigr|.
\]

\paragraph{Reference clip targets.}
\[
\epsilon_{\rm GSPO}^{(+)} \approx 3\!\times\!10^{-4},\qquad
\epsilon_{\rm GSPO}^{(-)} \approx 4\!\times\!10^{-4},\qquad
\epsilon_g^{\,\text{target}} \approx 3\!\times\!10^{-3}.
\]

% =============================================================================
\section{Trial plan}

Thirteen trials. Dimensions varied: \textsf{algo}
$\in\{\text{grpo}, \text{bspo}\}$, \textsf{allocation}
$\in\{\text{pro}, \text{const}, \text{ipro}\}$,
\textsf{adv} $\in\{\text{adv1}, \text{adv2}, \text{adv3}, \text{adv4}\}$,
\textsf{curriculum} $\in\{\text{on}, \text{off}\}$,
\textsf{penalty} $\in\{\text{on}(\lambda), \text{off}\}$,
\textsf{async\_clip} $\in\{\text{sync}, \text{async}\}$.

For \textsf{algo}=grpo (trials~1--2), BSPO is off: the loss is pure
verl \texttt{vanilla} (loss\_mode=vanilla) GRPO with the clip from
\texttt{actor\_rollout\_ref.actor.clip\_ratio\_\{low,high\}}. All
BSPO-specific dimensions (allocation, adv, penalty, async\_clip)
do not apply.

For \textsf{algo}=bspo (trials~3--13), the objective follows
\S2.7 with $V_{1a}$ when \textsf{penalty}=off, $V_{1b}$ when
\textsf{penalty}=on.

\paragraph{Trials, in priority-1 order.}

\begin{enumerate}[label=\textbf{Trial \arabic*}., wide=0pt]

\item \emph{(priority 1)} Pure GRPO baseline. No BSPO loss, no
  curriculum. Reference point for all BSPO vs GRPO comparisons.

\item \emph{(priority 2)} GRPO + curriculum. Isolates the
  contribution of the Gaussian curriculum sampler, independent of the
  BSPO loss.

\item \emph{(priority 3)} BSPO / \textsf{pro} / Adv~I / no curr /
  $V_{1a}$ / sync. The canonical BSPO configuration: proportional
  allocation, standard GRPO advantage (Adv~I by
  \eqref{eq:adv-I}), simplest objective. Baseline for all
  subsequent bspo ablations.

\item \emph{(priority 4)} Same as trial~3 plus curriculum. Isolates
  curriculum effect \emph{within} BSPO.

\item \emph{(priority 5)} Same as trial~3 plus $V_{1b}$ penalty with
  $\lambda = 1\!\times\!10^{-3}$. Tests whether the penalty term
  materially changes outcomes at low $\lambda$.

\item \emph{(priority 6)} Same as trial~5 with $\lambda =
  1\!\times\!10^{-2}$. Completes a one-decade $\lambda$ sweep.

\item \emph{(priority 7)} BSPO / \textsf{const} / Adv~I / no curr /
  $V_{1a}$ / sync. Swaps \textsf{pro} allocation for \textsf{const}
  (equal-per-group $\dbudget(g)$). const is the only allocation that
  forbids Adv~II/III/IV (Table~\ref{tab:adv-alloc}), so adv1 is
  mandatory.

\item \emph{(priority 8)} BSPO / \textsf{ipro} / Adv~I / no curr /
  $V_{1a}$ / sync. Swaps to \textsf{ipro} (inverse proportional).

\item \emph{(priority 9)} BSPO / \textsf{pro} / Adv~II / no curr /
  $V_{1a}$ / sync. Begins adv-variant sweep within pro allocation.

\item \emph{(priority 10)} BSPO / \textsf{pro} / Adv~III / no curr /
  $V_{1a}$ / sync. Reward-spread-only denominator.

\item \emph{(priority 11)} BSPO / \textsf{pro} / Adv~IV / no curr /
  $V_{1a}$ / sync. Reward-spread denominator scaled by $\hat d(g)$.

\item \emph{(priority 12)} BSPO / \textsf{ipro} / Adv~II / no curr /
  $V_{1a}$ / sync. The only non-pro adv-variant combination legal
  under ipro.

\item \emph{(priority 13)} BSPO / \textsf{pro} / Adv~I / no curr /
  $V_{1a}$ / \textbf{async}. Trial~3 with async clip. Tests whether
  the sign-split $\epsilon_\tau^{(\pm)}$ targeting the asymmetric
  $\epsilon_{\rm GSPO}$ reference ($3e\text{-}4$ / $4e\text{-}4$)
  changes the training dynamics. This is the only async-clip trial
  in the iter2 plan.

\end{enumerate}

Priority-1 is the intended execution order: baselines first, then the
pro-adv1 main axis (curriculum, penalty), then alternative
allocations, then alternative advantages, then async. The async trial
is last in queue because it is an orthogonal implementation detail
and its value depends on the baseline pro-adv1 result.

% =============================================================================
\appendix
\section{Trial table (raw)}

The trial plan in the machine-readable form used to schedule the sweep.
Blank cells in the CSV are rendered as ``---'' in
Table~\ref{tab:trials-raw}.

\begin{table}[h]
\centering
\caption{Iter2 trial plan.}
\label{tab:trials-raw}
\begin{tabular}{rrllllllr}
\toprule
trial\_id & priority & algo & allocation & adv & curr & penalty & $\lambda$ & async\_clip \\
\midrule
 1 &  1 & grpo & ---   & ---  & ---       & ---          & ---             & --- \\
 2 &  2 & grpo & ---   & ---  & with curr & ---          & ---             & --- \\
 3 &  3 & bspo & pro   & adv1 & ---       & ---          & ---             & --- \\
 4 &  4 & bspo & pro   & adv1 & with curr & ---          & ---             & --- \\
 5 &  5 & bspo & pro   & adv1 & ---       & with penalty & $1\!\times\!10^{-3}$ & --- \\
 6 &  6 & bspo & pro   & adv1 & ---       & with penalty & $1\!\times\!10^{-2}$ & --- \\
10 &  7 & bspo & const & adv1 & ---       & ---          & ---             & --- \\
11 &  8 & bspo & ipro  & adv1 & ---       & ---          & ---             & --- \\
 7 &  9 & bspo & pro   & adv2 & ---       & ---          & ---             & --- \\
 8 & 10 & bspo & pro   & adv3 & ---       & ---          & ---             & --- \\
 9 & 11 & bspo & pro   & adv4 & ---       & ---          & ---             & --- \\
12 & 12 & bspo & ipro  & adv2 & ---       & ---          & ---             & --- \\
13 & 13 & bspo & pro   & adv1 & ---       & ---          & ---             & async \\
\bottomrule
\end{tabular}
\end{table}

\noindent The rows are sorted here by priority so the appendix table
matches the execution order. Trial IDs 7--9 have priorities
9--11 because the adv-variant exploration is deferred until after
allocation sweeps (trials 10--11, which have priorities 7--8).

\end{document}
